\documentclass{ximera}

\input{../../../preamble.tex}


\definecolor{fvdnavy}{HTML}{071D43}
\definecolor{fvdcyan}{HTML}{20BFC6}
\definecolor{fvdcyanlight}{HTML}{EFFCFB}
\definecolor{fvdmagenta}{HTML}{F10D69}
\definecolor{fvdmagentalight}{HTML}{FFF1F7}
\definecolor{fvdtext}{HTML}{163044}





%=================================================
% Pedagogische omgevingen
% PDF: tcolorbox
% HTML: learning-box
%=================================================

\newenvironment{voorkennis}
{%
    \ifdefined\HCode
        \HCode{%
<section class="learning-box learning-box--prior">
<div class="learning-box__header">
<span class="learning-box__icon" aria-hidden="true"></span>
<span class="learning-box__title">Voorkennis</span>
</div>
<div class="learning-box__content">
        }%
        \begin{itemize}
    \else
       \begin{tcolorbox}[
    enhanced,
    breakable,
    colback=fvdcyanlight,
    colframe=fvdcyan,
    coltext=fvdtext,
    boxrule=0.5pt,
    leftrule=4pt,
    arc=4mm,
    outer arc=4mm,
    left=7mm,
    right=6mm,
    top=5mm,
    bottom=5mm,
    before skip=6mm,
    after skip=6mm,
    title=Voorkennis,
    coltitle=fvdcyan,
    fonttitle=\bfseries\Large,
    attach title to upper={\par\smallskip},
    boxed title style={
        empty,
        left=0mm,
        right=0mm,
        top=0mm,
        bottom=0mm
    },
    overlay={
        \draw[
            fvdcyan,
            line width=0.5pt
        ]
        ([xshift=42mm,yshift=-13mm]frame.north west)
        --
        ([xshift=-7mm,yshift=-13mm]frame.north east);
    }
]
        \begin{itemize}
    \fi
}
{%
    \end{itemize}
    \ifdefined\HCode
        \HCode{%
</div>
</section>
        }%
    \else
        \end{tcolorbox}
    \fi
}


\newenvironment{leerdoelen}
{%
    \ifdefined\HCode
        \HCode{%
<section class="learning-box learning-box--goals">
<div class="learning-box__header">
<span class="learning-box__icon" aria-hidden="true"></span>
<span class="learning-box__title">Leerdoelen</span>
</div>
<div class="learning-box__content">
        }%
        \begin{itemize}
    \else
\begin{tcolorbox}[
    enhanced,
    breakable,
    colback=fvdmagentalight,
    colframe=fvdmagenta,
    coltext=fvdtext,
    boxrule=0.5pt,
    leftrule=4pt,
    arc=4mm,
    outer arc=4mm,
    left=7mm,
    right=6mm,
    top=5mm,
    bottom=5mm,
    before skip=6mm,
    after skip=6mm,
    title=Leerdoelen,
    coltitle=fvdmagenta,
    fonttitle=\bfseries\Large,
    attach title to upper={\par\smallskip},
    boxed title style={
        empty,
        left=0mm,
        right=0mm,
        top=0mm,
        bottom=0mm
    },
    overlay={
        \draw[
            fvdmagenta,
            line width=0.5pt
        ]
        ([xshift=42mm,yshift=-13mm]frame.north west)
        --
        ([xshift=-7mm,yshift=-13mm]frame.north east);
    }
]
        \begin{itemize}
    \fi
}
{%
    \end{itemize}
    \ifdefined\HCode
        \HCode{%
</div>
</section>
        }%
    \else
        \end{tcolorbox}
    \fi
}


\addPrintStyle{../../..}

\title{Een beweging beschrijven}
\author{Ingmar Herreman}

\begin{abstract}
Beweging is overal: een fietser rijdt door het verkeer, een sprinter
vertrekt uit de startblokken, een trein rijdt een station binnen en
een satelliet beweegt rond de aarde.

In dit hoofdstuk bouwen we de taal op waarmee we zulke bewegingen
nauwkeurig kunnen beschrijven. We onderzoeken wat we bedoelen met
rust en beweging, kiezen een geschikt referentiestelsel en leren hoe
een werkelijk voorwerp kan worden vereenvoudigd tot een fysisch model.

Daarna beschrijven we de positie van een voorwerp en maken we
onderscheid tussen afgelegde weg en verplaatsing. We stellen
bewegingen voor met woorden, tabellen, grafieken en
functievoorschriften en onderzoeken hoe deze verschillende
voorstellingen met elkaar samenhangen.

Zo maken we de overgang van een waargenomen beweging naar een
wiskundig model. Dat vormt de basis voor de volgende hoofdstukken,
waarin we snelheid en versnelling gebruiken om bewegingen steeds
nauwkeuriger te analyseren.
\end{abstract}

\begin{document}

% FVD_CHAPTER_DATA_START
% Automatisch gegenereerd uit index.tex.
% Niet handmatig aanpassen.
\fvdchapterdata{1}{Beweging en kracht}{1}
% FVD_CHAPTER_DATA_END

\maketitle


% ============================================================
% VOORKENNIS EN LEERDOELEN
% ============================================================

\begin{voorkennis}
  \item Je kan een eenvoudige rechtlijnige beweging beschrijven.
  \item Je kan informatie aflezen uit een tabel en een grafiek.
  \item Je kent het verschil tussen een scalaire grootheid en een vector.
  \item Je kan met positieve en negatieve getallen werken.
  \item Je kan een eenvoudig functievoorschrift gebruiken om functiewaarden te berekenen.
\end{voorkennis}

\begin{leerdoelen}
  \item Je kan uitleggen waarom rust en beweging afhangen van het gekozen referentiestelsel.
  \item Je kan een beweging beschrijven met behulp van een referentiestelsel, positie en tijd.
  \item Je kan uitleggen wat het puntmassamodel inhoudt en beoordelen wanneer dat model bruikbaar is.
  \item Je kan baan, positie, afgelegde weg en verplaatsing van elkaar onderscheiden.
  \item Je kan bij een rechtlijnige beweging correct werken met positieve en negatieve posities en verplaatsingen.
  \item Je kan positie en verplaatsing vectorieel voorstellen en interpreteren.
  \item Je kan een beweging voorstellen met woorden, een tabel, een grafiek en een functievoorschrift.
  \item Je kan informatie uit een positie-tijddiagram aflezen en fysisch interpreteren.
  \item Je kan uitleggen waarom een meetreeks en een wiskundig model niet hetzelfde zijn.
  \item Je kan tussen verschillende voorstellingen van dezelfde beweging schakelen.
\end{leerdoelen}


% ============================================================
% 1 WAAROM?
% ============================================================

\FVDsection{Waarom leer je dit?}

Een auto die optrekt aan een verkeerslicht, een regendruppel die valt,
een voetbal die door de lucht vliegt en de aarde die rond de zon
beweegt: het zijn allemaal voorbeelden van beweging.

Op het eerste gezicht lijken deze bewegingen sterk van elkaar te
verschillen. Toch proberen we ze in de fysica met dezelfde begrippen
en wiskundige modellen te beschrijven.

De tak van de fysica die beweging en de oorzaken ervan bestudeert,
noemen we de \textbf{mechanica}.

Binnen de mechanica maken we een belangrijk onderscheid.

\begin{definitie}[Kinematica]

De \textbf{kinematica} beschrijft hoe een voorwerp beweegt zonder
de oorzaken van die beweging te onderzoeken.

\end{definitie}

\begin{definitie}[Dynamica]

De \textbf{dynamica} onderzoekt het verband tussen een beweging en
de krachten die deze beweging veroorzaken of veranderen.

\end{definitie}

Bij de kinematica stellen we bijvoorbeeld vragen als:

\begin{itemize}
  \item Waar bevindt een voorwerp zich?
  \item Welke baan volgt het?
  \item Hoe verandert zijn positie met de tijd?
  \item Hoe snel beweegt het?
  \item Verandert die snelheid?
\end{itemize}

Bij de dynamica stellen we uiteindelijk een andere vraag:

\begin{center}
\textit{Waarom beweegt het voorwerp op die manier?}
\end{center}

In dit hoofdstuk houden we ons voorlopig bezig met de eerste vraag.
We willen eerst leren hoe we een beweging nauwkeurig kunnen
\emph{beschrijven}.

\begin{opmerking}

De klassieke mechanica beschrijft zeer veel bewegingen uit onze
dagelijkse omgeving uitstekend. Ze is echter geen theorie die in
alle omstandigheden geldig is.

Later in deze cursus ontdekken we dat bij zeer hoge snelheden en op
zeer kleine schaal andere fysische modellen nodig zijn.

\end{opmerking}


% ============================================================
% 2 RUST EN BEWEGING
% ============================================================

\FVDsection{Rust en beweging}

Stel dat je in een trein zit die met constante snelheid door een
station rijdt. Op het tafeltje voor je ligt een boek.

Ten opzichte van het tafeltje verandert de positie van het boek
niet. We zeggen dat het boek \textbf{in rust is ten opzichte van de
trein}.

Voor iemand die op het perron staat, beweegt datzelfde boek echter
samen met de trein voorbij.

Hetzelfde voorwerp kan dus tegelijkertijd in rust zijn ten opzichte
van het ene voorwerp en bewegen ten opzichte van een ander voorwerp.

De uitspraak

\begin{center}
``Het boek beweegt.''
\end{center}

is daarom onvolledig.

We moeten vermelden \emph{ten opzichte waarvan} we de beweging
beschrijven.


\FVDsubsection{Het referentiestelsel}

\begin{definitie}[Referentiestelsel]

Een \textbf{referentiestelsel} is het stelsel ten opzichte waarvan
we de positie en de beweging van een voorwerp beschrijven.

\end{definitie}

Bij bewegingen dicht bij het aardoppervlak kiezen we vaak een
referentiestelsel dat vast verbonden is met de aarde.

\begin{eigenschap}[Relatief karakter van beweging]

Rust en beweging zijn relatieve begrippen.

Een voorwerp kan in rust zijn ten opzichte van het ene
referentiestelsel en tegelijk bewegen ten opzichte van een ander
referentiestelsel.

\end{eigenschap}


\begin{oefening}[Wie beweegt?]{\oefB\ESS}

Een leerling zit in een bus die langs een bushalte rijdt.

\begin{question}
Is de leerling in rust of in beweging ten opzichte van zijn stoel?
\end{question}

\begin{question}
Is de leerling in rust of in beweging ten opzichte van de bushalte?
\end{question}

\begin{question}
Is de buschauffeur in rust of in beweging ten opzichte van de leerling?
\end{question}

\begin{question}
Leg uit waarom de uitspraak
``de leerling beweegt''
zonder verdere informatie onvolledig is.
\end{question}

\end{oefening}


\begin{oefening}[Twee treinen]{\oefU}

Twee treinen rijden op evenwijdige sporen met dezelfde snelheid in
dezelfde richting.

Een reiziger in trein A kijkt naar een reiziger in trein B.

\begin{question}
Beschrijf de beweging van de reiziger in trein B ten opzichte van
trein A.
\end{question}

\begin{question}
Beschrijf dezelfde beweging ten opzichte van het station.
\end{question}

\begin{question}
Leg uit waarom beide beschrijvingen tegelijk correct kunnen zijn.
\end{question}

\end{oefening}


% ============================================================
% 3 MODEL
% ============================================================

\FVDsection{Van werkelijkheid naar model}

Werkelijke voorwerpen zijn vaak ingewikkeld.

Een auto heeft een bepaalde lengte en breedte. Zijn wielen draaien,
de motor bevat bewegende onderdelen en tijdens het rijden kunnen
verschillende delen van de auto verschillende bewegingen uitvoeren.

Wanneer we alleen willen onderzoeken hoe de auto zich langs een weg
verplaatst, hebben we al die details meestal niet nodig.

In de fysica vereenvoudigen we de werkelijkheid daarom met
\textbf{modellen}.

\[
\boxed{
\text{werkelijkheid}
\longrightarrow
\text{vereenvoudiging}
\longrightarrow
\text{fysisch model}
}
\]


\FVDsubsection{Het puntmassamodel}

Bij veel bewegingsproblemen stellen we een volledig voorwerp voor
door één enkel punt.

\begin{definitie}[Puntmassa]

Een \textbf{puntmassa} is een model waarbij we de afmetingen en de
vorm van een voorwerp verwaarlozen en de massa in één punt
geconcentreerd denken.

\end{definitie}

Dat betekent uiteraard niet dat het werkelijke voorwerp letterlijk
een punt is.

We kiezen het model omdat bepaalde eigenschappen van het voorwerp
voor onze onderzoeksvraag niet belangrijk zijn.

\begin{voorbeeld}[Een trein als puntmassa]

We willen bepalen hoe lang een trein nodig heeft om van Brussel naar
Gent te rijden.

Voor deze vraag kunnen we de trein meestal als puntmassa beschouwen.
Zijn lengte speelt nauwelijks een rol in het model.

Stellen we daarentegen de vraag wanneer de \emph{volledige} trein
een tunnel heeft verlaten, dan is de lengte van de trein wel
belangrijk.

Hetzelfde fysische systeem kan dus voor de ene onderzoeksvraag wel
en voor de andere niet als puntmassa worden gemodelleerd.

\end{voorbeeld}

\begin{opmerking}

Een fysisch model is een vereenvoudigde voorstelling van de
werkelijkheid.

De centrale vraag is niet of het model ``echt'' is, maar of de
vereenvoudigingen verantwoord zijn voor de onderzoeksvraag die we
willen beantwoorden.

\end{opmerking}


\begin{oefening}[Is een puntmassa geschikt?]{\oefB\ESS}

Geef aan of het puntmassamodel geschikt is. Motiveer telkens.

\begin{question}
We bepalen de tijd die een vliegtuig nodig heeft om van Brussel naar
New York te vliegen.
\end{question}

\begin{question}
We onderzoeken of een vrachtwagen onder een lage brug kan rijden.
\end{question}

\begin{question}
We bestuderen de baan van de aarde rond de zon.
\end{question}

\begin{question}
We bepalen hoe lang een sprinter over $100\,\mathrm{m}$ doet.
\end{question}

\end{oefening}


\begin{oefening}[Grenzen van een model]{\oefU\ESS}

Beschouw een satelliet die rond de aarde beweegt.

\begin{question}
Formuleer een onderzoeksvraag waarbij je de satelliet redelijk als
puntmassa kunt beschouwen.
\end{question}

\begin{question}
Formuleer een onderzoeksvraag waarbij het puntmassamodel onvoldoende is.
\end{question}

\begin{question}
Leg uit welke informatie in het tweede geval door het puntmassamodel
verloren gaat.
\end{question}

\end{oefening}


\begin{oefening}[De turner]{\oefG}

Een turner voert tijdens een sprong een salto uit.

Een leerling zegt:

\begin{quote}
``We kunnen de turner gewoon als een puntmassa beschouwen.''
\end{quote}

Een andere leerling zegt:

\begin{quote}
``Dat kan niet, want een turner heeft geen verwaarloosbare afmetingen.''
\end{quote}

Beoordeel beide uitspraken. Leg uit waarom het antwoord afhangt van
de onderzoeksvraag.

\end{oefening}


% ============================================================
% 4 BAAN
% ============================================================

\FVDsection{De baan van een beweging}

Tijdens een beweging neemt een puntmassa achtereenvolgens
verschillende posities in.

De verzameling van die opeenvolgende posities vormt de baan.

\begin{definitie}[Baan]

De \textbf{baan} of het \textbf{traject} van een bewegend voorwerp
is de verzameling posities die het voorwerp tijdens zijn beweging
inneemt.

\end{definitie}


\FVDsubsection{Rechtlijnige beweging}

Bij een \textbf{rechtlijnige beweging} is de baan een rechte.

Voorbeelden zijn:

\begin{itemize}
  \item een lift die verticaal beweegt;
  \item een trein op een recht stuk spoor;
  \item een voorwerp dat verticaal valt wanneer horizontale effecten
        verwaarloosbaar zijn.
\end{itemize}


\FVDsubsection{Kromlijnige beweging}

Bij een \textbf{kromlijnige beweging} is de baan een kromme.

Een geworpen basketbal volgt bijvoorbeeld een kromlijnige baan.


\FVDsubsection{Cirkelbeweging}

Een bijzonder geval van een kromlijnige beweging is de
\textbf{cirkelbeweging}. De baan is dan een cirkel.

Voorbeelden zijn:

\begin{itemize}
  \item een zitje van een reuzenrad;
  \item een punt op de rand van een draaiend fietswiel;
  \item de tip van een draaiende ventilator.
\end{itemize}

We beperken ons hier tot het herkennen van de baan. Later zullen we
de cirkelbeweging kwantitatief analyseren.


\begin{oefening}[Welke baan?]{\oefB}

Classificeer de baan als rechtlijnig, kromlijnig of cirkelvormig.

\begin{question}
Een lift beweegt van de benedenverdieping naar de zesde verdieping.
\end{question}

\begin{question}
Een voetbal wordt schuin omhoog geschoten.
\end{question}

\begin{question}
Een punt op de rand van een draaiende fietswielband.
\end{question}

\begin{question}
Een trein rijdt door een lang recht stuk spoor.
\end{question}

\end{oefening}


% ============================================================
% 5 POSITIE
% ============================================================

\FVDsection{Waar bevindt een voorwerp zich?}

Om een beweging nauwkeurig te beschrijven, moeten we kunnen aangeven
waar het bewegende voorwerp zich bevindt.

Voor een rechtlijnige beweging kiezen we een coördinaatas.

We leggen vast:

\begin{itemize}
  \item een oorsprong;
  \item een positieve richting;
  \item een eenheid.
\end{itemize}

Voor een beweging langs de $x$-as beschrijven we de positie met de
coördinaat $x$.

De positie op tijdstip $t$ noteren we als

\[
\boxed{x(t)}.
\]

De SI-eenheid van positie is meter:

\[
[x]=\mathrm{m}.
\]


\FVDsubsection{Een positie kan negatief zijn}

Stel dat we de positieve $x$-richting naar rechts kiezen.

Een voorwerp rechts van de oorsprong kan bijvoorbeeld de positie

\[
x=+5\,\mathrm{m}
\]

hebben.

Een voorwerp links van de oorsprong kan bijvoorbeeld de positie

\[
x=-3\,\mathrm{m}
\]

hebben.

Een negatieve positie is dus perfect mogelijk.

Het minteken vertelt aan welke kant van de oorsprong het voorwerp
zich bevindt volgens de gekozen positieve richting.

\begin{opmerking}

Een positie is geen afgelegde afstand.

De positie is een coördinaat binnen een gekozen referentiestelsel en
kan daarom positief, nul of negatief zijn.

\end{opmerking}


\begin{oefening}[Positie, weg of verplaatsing?]{\oefB\ESS}

Beoordeel de volgende uitspraken als mogelijk of onmogelijk.
Motiveer telkens.

\begin{question}
De positie van een voorwerp is $-30\,\mathrm{m}$.
\end{question}

\begin{question}
De afgelegde weg van een voorwerp is $-30\,\mathrm{m}$.
\end{question}

\begin{question}
De verplaatsing van een voorwerp is $-30\,\mathrm{m}$.
\end{question}

\end{oefening}


% ============================================================
% 6 WEG EN VERPLAATSING
% ============================================================

\FVDsection{Afgelegde weg en verplaatsing}

De begrippen \textbf{afgelegde weg} en \textbf{verplaatsing} worden
in het dagelijks taalgebruik soms door elkaar gebruikt. In de fysica
betekenen ze iets verschillends.

Beschouw een persoon die start op

\[
x_A=10\,\mathrm{m},
\]

vervolgens naar

\[
x_B=50\,\mathrm{m}
\]

loopt en daarna terugkeert naar

\[
x_C=30\,\mathrm{m}.
\]


\FVDsubsection{Afgelegde weg}

Tijdens het eerste deel van de beweging legt de persoon af:

\[
50-10=40\,\mathrm{m}.
\]

Tijdens het tweede deel:

\[
50-30=20\,\mathrm{m}.
\]

De totale afgelegde weg bedraagt dus:

\[
\Delta s=40+20=60\,\mathrm{m}.
\]

\begin{definitie}[Afgelegde weg]

De \textbf{afgelegde weg} is de totale lengte van het traject dat
een voorwerp tijdens een beweging doorloopt.

\end{definitie}

De afgelegde weg is niet negatief:

\[
\boxed{\Delta s\geq0}.
\]


\FVDsubsection{Verplaatsing}

Voor de verplaatsing kijken we alleen naar de begin- en eindpositie.

\[
\Delta x=x_C-x_A.
\]

Dus:

\[
\Delta x=30-10=20\,\mathrm{m}.
\]

\begin{definitie}[Verplaatsing]

De \textbf{verplaatsing} bij een rechtlijnige beweging is het
verschil tussen de eindpositie en de beginpositie:

\[
\boxed{
\Delta x=x_{\mathrm{eind}}-x_{\mathrm{begin}}.
}
\]

\end{definitie}

Een verplaatsing kan positief, nul of negatief zijn.


\FVDsubsection{Vergelijk beide grootheden}

In ons voorbeeld vonden we:

\[
\Delta s=60\,\mathrm{m}
\]

maar:

\[
\Delta x=20\,\mathrm{m}.
\]

De afgelegde weg houdt rekening met het volledige traject.
De verplaatsing vergelijkt alleen begin- en eindpositie.

Daarom geldt bij een rechtlijnige beweging:

\[
\boxed{
|\Delta x|\leq\Delta s.
}
\]

De gelijkheid geldt wanneer het voorwerp tijdens het beschouwde
tijdinterval niet van bewegingsrichting verandert.

\begin{voorbeeld}[Een ronde op de atletiekpiste]

Een hardloper loopt één volledige ronde van $400\,\mathrm{m}$ en
eindigt opnieuw op zijn startplaats.

Dan is

\[
\Delta s=400\,\mathrm{m},
\]

maar

\[
\Delta x=0.
\]

De hardloper heeft dus wel degelijk bewogen, hoewel zijn verplaatsing
nul is.

\end{voorbeeld}


\begin{oefening}[Een robot op een rechte]{\oefB\ESS}

Een robot start op $x=-2,0\,\mathrm{m}$, rijdt naar
$x=7,0\,\mathrm{m}$ en keert daarna terug naar
$x=3,0\,\mathrm{m}$.

\begin{question}
Bepaal de eindpositie.
\end{question}

\begin{question}
Bepaal de totale afgelegde weg.
\end{question}

\begin{question}
Bepaal de verplaatsing.
\end{question}

\end{oefening}


\begin{oefening}[Nul verplaatsing]{\oefU\ESS}

Een leerling beweert:

\begin{quote}
``Als de verplaatsing nul is, heeft het voorwerp niet bewogen.''
\end{quote}

\begin{question}
Geef een concreet tegenvoorbeeld.
\end{question}

\begin{question}
Leg uit waarom de uitspraak fout is.
\end{question}

\end{oefening}


\begin{oefening}[Kan de verplaatsing groter zijn?]{\oefU}

Kan de grootte van de verplaatsing ooit groter zijn dan de afgelegde
weg?

Formuleer een argument voor je antwoord.

\end{oefening}


% ============================================================
% 7 VECTOR
% ============================================================

\FVDsection{Positie en verplaatsing als vector}

Voor een beweging langs één rechte volstaat één coördinaat $x$.
Bewegingen in de werkelijkheid zijn echter vaak twee- of
driedimensionaal.

Daarom zullen we bewegingen later vectorieel beschrijven.

\begin{definitie}[Positievector]

De \textbf{positievector}

\[
\vec r(t)
\]

loopt van de oorsprong van het gekozen coördinatenstelsel naar de
positie van het voorwerp op tijdstip $t$.

\end{definitie}

In twee dimensies kunnen we schrijven:

\[
\boxed{
\vec r(t)=x(t)\vec e_x+y(t)\vec e_y.
}
\]

Hierin zijn $x(t)$ en $y(t)$ de componenten van de positievector.


\FVDsubsection{Verplaatsingsvector}

De verplaatsing tussen twee tijdstippen $t_1$ en $t_2$ is:

\[
\boxed{
\Delta\vec r
=
\vec r(t_2)-\vec r(t_1).
}
\]

Voor een beweging langs de $x$-as reduceert dit tot:

\[
\Delta x=x(t_2)-x(t_1).
\]

\begin{opmerking}

De vectoriële beschrijving lijkt voor een zuiver rechtlijnige
beweging misschien omslachtig. Ze wordt echter onmisbaar wanneer we
later bewegingen in twee dimensies onderzoeken.

Daar zullen ook de snelheids- en versnellingsvector een centrale rol
spelen.

\end{opmerking}


\begin{oefening}[Verplaatsingsvector]{\oefB}

Een drone bevindt zich achtereenvolgens op

\[
\vec r_1=
\begin{pmatrix}
2\\
3
\end{pmatrix}\mathrm{m}
\]

en

\[
\vec r_2=
\begin{pmatrix}
7\\
1
\end{pmatrix}\mathrm{m}.
\]

Bepaal

\[
\Delta\vec r=\vec r_2-\vec r_1.
\]

\end{oefening}


\begin{oefening}[Traject of verplaatsing?]{\oefU\ESS}

Een wandelaar beweegt in een vlak van punt $A$ naar punt $B$.

\begin{question}
Leg uit waarom de verplaatsingsvector alleen afhangt van $A$ en $B$.
\end{question}

\begin{question}
Leg uit waarom de afgelegde weg wel afhangt van het gevolgde traject.
\end{question}

\end{oefening}


% ============================================================
% 8 TIJD
% ============================================================

\FVDsection{Beweging betekent verandering in de tijd}

Een positie alleen is onvoldoende om een beweging te beschrijven.

Stel dat twee auto's dezelfde verplaatsing maken. De ene auto doet
daar $5\,\mathrm{s}$ over en de andere $20\,\mathrm{s}$.

De begin- en eindpositie kunnen dezelfde zijn, maar de beweging is
duidelijk niet dezelfde.

We moeten daarom ook vastleggen \emph{wanneer} een voorwerp een
bepaalde positie inneemt.

De positie wordt een functie van de tijd:

\[
\boxed{x=x(t)}
\]

of korter:

\[
\boxed{x(t)}.
\]

De onafhankelijke veranderlijke is de tijd $t$.
De afhankelijke veranderlijke is de positie $x$.


% ============================================================
% 9 REPRESENTATIES
% ============================================================

\FVDsection{Een beweging op verschillende manieren voorstellen}

Eenzelfde beweging kan op verschillende manieren worden beschreven.

\[
\boxed{
\text{situatie}
\longleftrightarrow
\text{tabel}
\longleftrightarrow
\text{grafiek}
\longleftrightarrow
\text{functievoorschrift}
}
\]

Deze voorstellingen beschrijven dezelfde fysische beweging, maar
maken verschillende eigenschappen zichtbaar.


\FVDsubsection{Een beweging in een tabel}

Beschouw de volgende meetgegevens:

\[
\begin{array}{c|ccccc}
t\;(\mathrm{s}) & 0 & 1 & 2 & 3 & 4\\ \hline
x\;(\mathrm{m}) & 5 & 8 & 11 & 14 & 17
\end{array}
\]

Uit de tabel lezen we bijvoorbeeld af dat het voorwerp op

\[
t=2\,\mathrm{s}
\]

de positie

\[
x=11\,\mathrm{m}
\]

inneemt.


\FVDsubsection{Een positie-tijddiagram}

Dezelfde meetgegevens kunnen we voorstellen in een grafiek.

Op de horizontale as plaatsen we de tijd $t$ en op de verticale as
de positie $x$.

Zo ontstaat een \textbf{positie-tijddiagram} of een
\textbf{$x(t)$-diagram}.

\begin{opmerking}

Een $x(t)$-grafiek is niet de baan van het voorwerp.

De horizontale as stelt de tijd voor, niet een tweede
plaatscoördinaat.

\end{opmerking}

Een rechte of kromme in een positie-tijddiagram moet dus worden
geïnterpreteerd als een verband tussen positie en tijd.


\begin{oefening}[Grafiek of baan?]{\oefU\ESS}

Een leerling bekijkt een $x(t)$-grafiek die een parabool vormt en zegt:

\begin{quote}
``Het voorwerp beweegt dus langs een parabolische baan.''
\end{quote}

\begin{question}
Is die conclusie correct?
\end{question}

\begin{question}
Leg nauwkeurig uit wat de grafiek wel en niet voorstelt.
\end{question}

\end{oefening}


% ============================================================
% 10 PLAATSFUNCTIE
% ============================================================

\FVDsection{De plaatsfunctie}

Wanneer we het verband tussen positie en tijd met een functie kunnen
beschrijven, krijgen we een wiskundig model van de beweging.

\begin{definitie}[Plaatsfunctie]

Een \textbf{plaatsfunctie} of \textbf{bewegingsvergelijking}
beschrijft de positie van een voorwerp als functie van de tijd.

Voor een rechtlijnige beweging schrijven we:

\[
\boxed{x=x(t).}
\]

\end{definitie}


\begin{voorbeeld}[Een eenvoudige plaatsfunctie]

Beschouw:

\[
x(t)=2,0+3,0t,
\]

waarbij $x$ in meter en $t$ in seconde wordt uitgedrukt.

Op het begintijdstip:

\[
x(0)=2,0\,\mathrm{m}.
\]

Na $4,0\,\mathrm{s}$:

\[
x(4,0)
=
2,0+3,0\cdot4,0
=
14,0\,\mathrm{m}.
\]

Het getal $2,0$ vertelt ons dus de beginpositie.

De factor $3,0$ zegt iets over hoe snel de positie verandert.
De precieze fysische betekenis daarvan onderzoeken we in het
volgende hoofdstuk.

\end{voorbeeld}


\begin{oefening}[Werken met een plaatsfunctie]{\oefB\ESS}

Een bewegend voorwerp wordt beschreven door

\[
x(t)=5,0-2,0t,
\]

met $x$ in meter en $t$ in seconde.

\begin{question}
Bepaal $x(0)$.
\end{question}

\begin{question}
Bepaal $x(2,0)$.
\end{question}

\begin{question}
Bepaal $x(4,0)$.
\end{question}

\begin{question}
Bepaal de verplaatsing tussen $t=1,0\,\mathrm{s}$ en
$t=4,0\,\mathrm{s}$.
\end{question}

\begin{question}
Beschrijf in woorden wat je aan de waarden van $x(t)$ opmerkt.
\end{question}

\end{oefening}


\begin{oefening}[Van tabel naar model]{\oefU\ESS}

Van een beweging zijn de volgende gegevens bekend:

\[
\begin{array}{c|ccccc}
t\;(\mathrm{s}) & 0 & 2 & 4 & 6 & 8\\ \hline
x\;(\mathrm{m}) & -3 & 1 & 5 & 9 & 13
\end{array}
\]

\begin{question}
Stel de gegevens grafisch voor.
\end{question}

\begin{question}
Zoek een functievoorschrift dat bij de gegevens past.
\end{question}

\begin{question}
Welke fysische betekenis heeft de constante term?
\end{question}

\begin{question}
Gebruik je model om de positie op $t=5\,\mathrm{s}$ te voorspellen.
\end{question}

\end{oefening}


% ============================================================
% 11 MODELLEREN
% ============================================================

\FVDsection{Van meetgegevens naar een model}

In een echt experiment krijgen we meestal niet onmiddellijk een
functievoorschrift.

We meten eerst waarden:

\[
(t_1,x_1),\quad
(t_2,x_2),\quad
(t_3,x_3),\quad\ldots
\]

Die meetgegevens kunnen we in een tabel plaatsen en vervolgens
grafisch voorstellen.

Daarna zoeken we een functie die het waargenomen verband zo goed
mogelijk beschrijft.

\[
\boxed{
\text{werkelijke beweging}
\longrightarrow
\text{metingen}
\longrightarrow
\text{grafiek}
\longrightarrow
\text{wiskundig model}
}
\]

\begin{opmerking}

Meetgegevens en een wiskundig model zijn niet hetzelfde.

Meetpunten zijn experimentele waarnemingen. Een functievoorschrift
is een vereenvoudigd model waarmee we proberen het patroon in die
gegevens te beschrijven.

\end{opmerking}


\FVDsubsection{Voorbeeld: een sprint analyseren}

Bij een sprint kunnen tussentijden worden gemeten op verschillende
posities.

We verzamelen bijvoorbeeld gegevens van de vorm:

\[
\begin{array}{c|cccccc}
t\;(\mathrm{s}) & t_0&t_1&t_2&t_3&t_4&t_5\\ \hline
x\;(\mathrm{m}) & 0&20&40&60&80&100
\end{array}
\]

en stellen de punten $(t,x)$ grafisch voor.

Met een digitaal hulpmiddel kunnen we vervolgens onderzoeken welk
type functie de meetgegevens voldoende goed beschrijft.

Het gevonden functievoorschrift is geen exacte ``wet'' voor de
sprinter. Het is een model van de gemeten beweging binnen het
onderzochte tijdsinterval.


\begin{oefening}[Een beweging modelleren met ICT]{\oefU\ESS}

Je krijgt een tabel met gemeten posities van een bewegend voorwerp.

\begin{question}
Maak met een grafisch hulpmiddel een spreidingsdiagram van $x$ als
functie van $t$.
\end{question}

\begin{question}
Beschrijf het patroon in de meetpunten.
\end{question}

\begin{question}
Zoek een geschikt wiskundig model.
\end{question}

\begin{question}
Vergelijk enkele waarden die het model voorspelt met de oorspronkelijke
meetgegevens.
\end{question}

\begin{question}
Geef minstens één reden waarom een wiskundig model nooit volledig
samenvalt met de werkelijke beweging.
\end{question}

\end{oefening}


% ============================================================
% 12 INTERPRETEREN
% ============================================================

\FVDsection{Een bewegingsgrafiek interpreteren}

Een grafiek is meer dan een afbeelding van een formule.

Uit een positie-tijddiagram kunnen we rechtstreeks fysische
informatie halen.

We kunnen bijvoorbeeld onderzoeken:

\begin{itemize}
  \item waar een voorwerp zich op een bepaald tijdstip bevindt;
  \item wanneer het voorwerp de oorsprong passeert;
  \item wanneer het zich links of rechts van de oorsprong bevindt;
  \item wat de verplaatsing tussen twee tijdstippen is;
  \item of het voorwerp op verschillende tijdstippen dezelfde
        positie inneemt.
\end{itemize}


\begin{oefening}[Een plaatsfunctie onderzoeken]{\oefU\ESS}

Voor een bewegend voorwerp geldt:

\[
x(t)
=
-1,20t^3
+4,50t^2
-4,00t
+2,00,
\]

waarbij $x$ in meter en $t$ in seconde wordt uitgedrukt.

Onderzoek de beweging voor $t\geq0$.

\begin{question}
Maak een grafiek van $x(t)$ met een geschikt digitaal hulpmiddel.
\end{question}

\begin{question}
Bepaal de beginpositie.
\end{question}

\begin{question}
Bepaal de positie op $t=1,0\,\mathrm{s}$ en $t=2,0\,\mathrm{s}$.
\end{question}

\begin{question}
Onderzoek wanneer het voorwerp de oorsprong passeert.
\end{question}

\begin{question}
Bepaal de verplaatsing tussen $t=0$ en $t=2,0\,\mathrm{s}$.
\end{question}

\begin{question}
Zoek, indien mogelijk, twee verschillende tijdstippen waarop het
voorwerp dezelfde positie inneemt.
\end{question}

\begin{question}
Leg uit waarom de positie alleen nog niet volledig vertelt hoe snel
het voorwerp beweegt.
\end{question}

\end{oefening}


\begin{oefening}[Ontwerp zelf een beweging]{\oefG}

Ontwerp zelf een plaatsfunctie $x(t)$ van een voorwerp dat:

\begin{itemize}
  \item start rechts van de oorsprong;
  \item op een later tijdstip de oorsprong passeert;
  \item daarna opnieuw terugkeert naar zijn beginpositie.
\end{itemize}

\begin{question}
Geef een mogelijk functievoorschrift.
\end{question}

\begin{question}
Maak een grafiek van je model.
\end{question}

\begin{question}
Toon aan dat je model aan de drie voorwaarden voldoet.
\end{question}

\begin{question}
Zijn er verschillende antwoorden mogelijk? Verklaar.
\end{question}

\end{oefening}


% ============================================================
% 13 SAMENVATTING
% ============================================================

\FVDsection{Wat hebben we opgebouwd?}

We begonnen dit hoofdstuk met een werkelijke beweging.

Om die beweging fysisch te beschrijven, hebben we stap voor stap een
model opgebouwd:

\[
\boxed{
\text{werkelijke beweging}
\longrightarrow
\text{puntmassa}
\longrightarrow
\text{referentiestelsel}
\longrightarrow
\text{positie}
\longrightarrow
x(t)
}
\]

We leerden dat rust en beweging afhangen van het gekozen
referentiestelsel.

We maakten onderscheid tussen de afgelegde weg

\[
\Delta s
\]

en de verplaatsing

\[
\Delta x=x_{\mathrm{eind}}-x_{\mathrm{begin}}.
\]

Voor bewegingen in meerdere dimensies introduceerden we de
positievector

\[
\vec r(t)
\]

en de verplaatsingsvector

\[
\Delta\vec r.
\]

Ten slotte beschreven we een beweging met tabellen, grafieken en
functievoorschriften:

\[
\boxed{
\text{tabel}
\longleftrightarrow
\text{grafiek}
\longleftrightarrow
\text{functie}
}
\]

Deze verschillende voorstellingen vormen samen een krachtige taal om
bewegingen te beschrijven en te modelleren.


% ============================================================
% 14 VOORUITBLIK
% ============================================================

\FVDsection{Vooruitblik: hoe snel verandert de positie?}

Stel dat twee auto's vertrekken op dezelfde positie en na een rit
dezelfde eindpositie bereiken.

De eerste auto doet daar $10\,\mathrm{s}$ over en de tweede
$25\,\mathrm{s}$.

Hun verplaatsing kan dezelfde zijn, maar hun beweging is duidelijk
niet dezelfde.

Alleen weten \emph{waar} een voorwerp zich bevindt, is dus niet
voldoende.

We willen ook weten:

\begin{center}
\textbf{Hoe snel verandert de positie?}
\end{center}

Wiskundig betekent dit dat we willen onderzoeken hoe de functie
$x(t)$ verandert wanneer de tijd verandert.

Dat brengt ons in het volgende hoofdstuk bij een nieuwe fysische
grootheid:

\[
\boxed{\text{snelheid}}
\]



\end{document}